% SNN 信息丢失与恢复：中文汇总与研究计划
\documentclass[journal,12pt,onecolumn]{IEEEtran}
\IEEEoverridecommandlockouts
\usepackage{cite}

\usepackage{amsmath,amssymb,amsfonts}
\usepackage{algorithmic}
\usepackage{algorithm}
%\usepackage{algpseudocode}
\usepackage{graphicx}

\usepackage{booktabs} 

\usepackage{textcomp}
%\usepackage[marginal]{footmisc}
\usepackage{xcolor}
\usepackage{subfigure}
\graphicspath{{./img/}}
\usepackage{tabularx}
\usepackage{makecell}
\usepackage{caption2}
%\usepackage{fontspec}
%\setmainfont{TeX Gyre Termes}
\usepackage{color}

\usepackage[UTF8]{ctex}

\begin{document}
\title{脉冲神经网络（SNN）训练中的信息丢失与恢复：中文汇总与研究思路}
\author{Tao}
\date{\today}
\maketitle

\begin{abstract}
本文基于仓库中 `Algorithm/` 目录下的 8 份 PDF 文献，采用 Python 程序将其转化为纯文本，并在 \texttt{Work\_Copilot/snn\_info\_summary\_cn.tex} 中完成中文汇总。本文包括：(1) 现有 SNN 算法的主要机制与信息丢失环节；(2) 常用“模糊边界/近似”技术的作用机理；(3) 可量化指标与恢复机制；(4) 基于当前证据的研究思路与后续计划。
\end{abstract}

\section{高层概要（概括性结论）}
\begin{itemize}
  \item SNN 训练中信息丢失主要来自离散化（阈值/重置）、时间截断（窗口化/衰减）、局部化近似（资格迹或随机反馈）、以及替代梯度导致的梯度偏差。\
  \item 常用缓解手段包括替代梯度（surrogate）、资格迹（三因子规则）、随机/固定反馈近似、以及 ANN→SNN 的软化映射；这些方法通常以牺牲精确度为代价换取可训练性与在线/低内存实现。\
  \item 信息丢失可以用互信息、重建误差、Spike-timing 距离、Fisher 信息近似等多维指标量化；区分"压缩"与"不可逆丢失"需要基于最优解码器的可恢复性测试。\
  \item 恢复策略可发生在训练内（资格迹、动态阈值、潜变量推断）、训练后（强解码器、时间重采样、生成式重建）或利用约束反向估计（把截断梯度视为 Fisher 近似用于校正）。
\end{itemize}

\section{方法族与仓库覆盖}
\begin{enumerate}
  \item BPTT（完整反向传播通过时间）：基线，能精确分配时间信用，但需 $O(T)$ 内存且存在梯度消失/爆炸问题。
  \item 三因子/资格迹类（S‑TLLR、TESS）：局部在线规则，使用资格迹 $e_{ij}[t]$ 与学习信号 $\delta_i[t]$ 的乘积更新权重，适合在线学习和低内存场景。\
  \item 事件驱动/在线近似类（OTTT、NDOT）：本仓库包含这两类方法的论文，侧重前向时间学习与常数训练记忆 \cite{jiang2024ndot,xiao2022online}.\
  \item 替代梯度家族（surrogate gradient）：通过平滑不可导阶跃激活恢复梯度信息流；NDOT/OTTT 文本中均提及该技术 \cite{xiao2022online}.\
  \item ANN→SNN 转换：率或延迟编码替代连续激活，通常丢失单脉冲时序细节，为对比基线提供分析视角。
\end{enumerate}

\section{导致信息丢失的具体环节（分类与机理）}
\subsection{时间维度}
\begin{itemize}
  \item 时间量化与采样间隔：采样步长 dt 决定可表示的频带，粗采样会丢失高频时间信息。\
  \item 资格迹的衰减：以指数衰减形式近似长期依赖，导致远期事件对当前学习信号的影响被模糊化。\
  \item 截断 BPTT 或在线近似：当仅利用有限历史或即时损失时，超出窗口的时间贡献被舍弃。
\end{itemize}

\subsection{空间/表示维度}
\begin{itemize}
  \item 局部更新：三因子规则和本地学习规则仅利用局部可用信号，无法完整表达跨层全局梯度结构。\
  \item 量化/稀疏化：为节省带宽和能耗，对膜电位、权重或脉冲进行量化，会丢失连续幅值信息。\
\end{itemize}

\subsection{梯度与优化维度}
\begin{itemize}
  \item 替代梯度的失真：用连续代理函数近似阶跃函数导数时，会改变梯度幅值和二阶特性，但保留方向信息的主体。\
  \item 近似反馈：如随机反馈或反馈对齐用近似梯度方向替代精确反向传播，会引入方向噪声。\
\end{itemize}

\subsection{离散非线性截断}
\begin{itemize}
  \item 硬阈值和完全重置使膜电位信息在放电后不可逆丢失，这类似于对连续信号的非可逆量化。\
\end{itemize}

\section{“模糊边界/近似”技术与信息恢复机理}
\begin{enumerate}[leftmargin=*]
  \item 替代梯度（surrogate gradient）
  \begin{itemize}
    \item 机理：用可导函数替代 Heaviside 的导数，使误差能够传播回脉冲生成模块。\
    \item 信息影响：恢复梯度方向信息，减少不可导引起的零梯度问题；但幅值和高阶信息仍可能失真。\
  \end{itemize}
  \item 资格迹与三因子规则
  \begin{itemize}
    \item 机理：利用递推变量记录前后脉冲相关性，并用学习信号调制，保留历史时序统计。\
    \item 信息影响：在有限记忆下保留时序相关，但对远期或高频信息存在衰减。\
  \end{itemize}
  \item 本地学习/反馈近似
  \begin{itemize}
    \item 机理：以固定随机矩阵或局部学习信号替代层间反向传播，降低空间复杂度。\
    \item 信息影响：保留误差信号的粗粒度流动，可训练性增强，但精确梯度空间结构部分丢失。\
  \end{itemize}
  \item ANN→SNN 软化映射
  \begin{itemize}
    \item 机理：通过平均率或时间延迟将连续激活映射到脉冲，以平滑离散化误差。\
    \item 信息影响：保留整体信息量，但丢失单脉冲级别时间细节。\
  \end{itemize}
\end{enumerate}

\section{可量化指标建议}
\begin{itemize}
  \item 互信息 $I(X;S)$：输入与脉冲序列之间的信息上限，反映可恢复性。\
  \item Fisher 信息与其近似：衡量参数变化对输出的敏感度，可用于评估梯度近似带来的信息损失。\
  \item 重建误差（MSE/SNR）：训练解码器从 spikes 恢复原始信号，直接衡量实际可恢复信息。\
  \item Spike-timing 距离：如 Victor--Purpura、van Rossum，用于时序抖动的量化。\
  \item 活动统计：发放率、脉冲覆盖度、熵、有效秩，衡量表征空间的富集程度。\
  \item 任务性能：分类/回归精度作为间接信息保留指标。\
\end{itemize}

\section{信息恢复机制（训练内/训练后/约束反向）}
\subsection{训练阶段内恢复}
\begin{itemize}
  \item 资格迹和三因子规则（S-TLLR、TESS）：在线记录时序相关，作为梯度近似的记忆载体。\
  \item 动态阈值与软重置：降低硬阈值的不可逆效应，保留膜电位残余信息。\
  \item 潜在变量推断：将膜电位视为隐变量，通过后验推断补偿观测脉冲带来的信息损失。\
  \item 增强 surrogate 设计：用更精细的近似函数保持更多梯度细节。\
\end{itemize}

\subsection{训练阶段后恢复}
\begin{itemize}
  \item 强解码器/群体解码：用并行神经元的冗余信息恢复细节。\
  \item 时间重采样和插值：基于先验在离散脉冲间补偿高频成分。\
  \item 生成式恢复：用生成模型或扩散模型从脉冲恢复连续信号。\
  \item 多脉冲/多位宽扩展：通过增加脉冲载体提高信息容量。\
\end{itemize}

\subsection{约束反向恢复}
\begin{itemize}
  \item 将截断梯度或替代梯度视为 Fisher 信息的粗估，用于优化器的自适应预条件或不确定性估计。\
  \item 变分框架：直接在观测与隐变量间推断，重建被丢弃的连续状态。\
\end{itemize}

\section{Algorithm 论文汇总与证据}
\subsection{转换结果}
已将 \texttt{Algorithm/} 下 8 个 PDF 文献转换为文本，保存在 \texttt{Algorithm/converted\_texts/} 下，便于逐页检索与内容摘录。\\

\subsection{NDOT}
\begin{itemize}
  \item 论文名：`NDOT: Neuronal Dynamics-based Online Training for Spiking Neural Networks`，并包含一个 bilinear 版本 \cite{jiang2024ndot}.\
  \item 关键点：全梯度分解为时间梯度和空间梯度，使用神经元动力学连续信息实现前向时间训练；强调在线性和常数内存开销。\
  \item 术语：包含 surrogate gradient、online training、temporal gradient、spatial gradient；未直接使用“信息丢失”。\
\end{itemize}

\subsection{OTTT}
\begin{itemize}
  \item 论文名：`Online Training Through Time for Spiking Neural Networks` \cite{xiao2022online}.\
  \item 关键点：通过跟踪 presynaptic activity 及瞬时损失，实现前向时间学习；证明 OTTT 与 spike representation 梯度具有类似下降方向，并提出三因子 Hebbian 更新形式。\
  \item 术语：强调 online learning、surrogate gradient、equivalent mapping、three-factor Hebbian；未直接使用“信息丢失”。\
\end{itemize}

\subsection{S-TLLR}
\begin{itemize}
  \item 论文名：`S-TLLR: STDP-inspired Temporal Local Learning Rule for Spiking Neural Networks` \cite{apolinario2023s}.\
  \item 关键点：STDP 灵感的时间局部三因子学习规则，兼顾因果与非因果时序关系，内存复杂度与时间步数无关。\
  \item 术语：强调 local learning、eligibility traces、temporal locality；当前文本未直接提及“信息丢失”。\
\end{itemize}

\subsection{TESS}
\begin{itemize}
  \item 论文名：`TESS: A Scalable Temporally and Spatially Local Learning Rule for Spiking Neural Networks` \cite{apolinario2025tess}.\
  \item 关键点：实现时间局部与空间局部，使用本地信号和同步学习信号，避免跨层误差反向传播；内存复杂度与时间步数无关。\
  \item 术语：强调 temporally local、spatially local、eligibility traces、local learning signal；未直接使用“信息丢失”。\
\end{itemize}

\subsection{当前证据与缺失}
\begin{itemize}
  \item 当前转换文本中可直接确认的主要内容：online training、surrogate gradient、eligibility trace、local learning、constant memory、three-factor update 等。\
  \item 但从文本检索结果来看，论文中尚未直接明确出现“information loss”或“信息丢失”的术语。\
  \item 因此，本报告的“信息丢失与恢复”分析需进一步结合具体段落或作者原话进行补充验证。\
\end{itemize}

\section{研究思路与后续计划}
\subsection{候选算法族}
\begin{itemize}
  \item 替代梯度类型：BPTT+SG 基线与其在线改进（NDOT、OTTT）。\
  \item 事件驱动在线局部类型：S-TLLR、TESS，关注资格迹和本地学习信号的时间/空间信息保留。\
  \item ANN→SNN 类型：用于对比离散化与时间信息丢失。\
\end{itemize}

\subsection{建议实验设计}
\begin{itemize}
  \item 任务：输入噪声控制的信号重建、事件流分类、CIFAR10-DVS、IBM DVS Gesture。\
  \item 变量：时间步长、噪声强度、eligibility decay、surrogate 形状、在线/离线训练、网络冗余度。\
  \item 指标：重建 MSE/SNR、互信息估计、Spike-timing 距离、脉冲活动覆盖度、分类精度。\
  \item 区分压缩与丢失：若强解码器能恢复信号则更接近压缩；若不可恢复且互信息下降则判定为信息丢失。\
\end{itemize}

\subsection{可能的新方向}
\begin{itemize}
  \item 变分推理：把膜电位建模为隐变量，利用后验恢复丢弃的连续内容。\
  \item 扩散/生成式模型：从脉冲序列恢复高频时间信息。\
  \item 记忆重放与历史事件回放：补偿在线近似带来的长期依赖缺失。\
\end{itemize}

\section*{结论}
本次工作已完成：
\begin{itemize}
  \item 在 \texttt{Algorithm/} 下建立 Python 工作区，生成 \texttt{convert\_pdfs.py}.\\
  \item 将 8 份 PDF 转换为文本，保存在 \texttt{Algorithm/converted\_texts/}。\\
  \item 在本 LaTeX 文档中按 `NDOT`、`OTTT`、`S-TLLR`、`TESS` 进行中文汇总，并注明当前证据与缺失。\
\end{itemize}

后续建议：请继续提供具体摘要或关键段落，将根据这些原文进一步补充“信息丢失”、“恢复机制”和“定量指标”的精确引用。

\bibliographystyle{IEEEtran}
\bibliography{../SNNInformationReconstruction}
\end{document}
